\section{讨论与结论}

本节用于控制论文边界并回到核心故事。结论不应展开所有实现细节，而应强调：\SPMPC{} 将液体晃动视为在线局部规划时域内的预测动态状态，而不是事后的平滑性指标。

\subsection{讨论边界}

第一，低阶晃液模型是面向实时优化的近似。它能够表示主导晃液模态和历史激励的动态记忆，但不能覆盖所有非线性自由液面现象，例如破波、冲击、强三维效应和复杂容器几何。

第二，内部模型 proxy 与真实液面评价必须区分。\texttt{/spmpc/slosh\_height} 或其它模型预测量用于诊断规划行为和解释 OCP 决策；真实实验结论应主要依赖外部液面观测，例如 RGB max-LCR、液面轮廓变化或其它独立传感指标。

第三，本文方法不提供形式化 spill-free guarantee。即使使用可选的 model-predicted slosh bound，也只能说明规划器限制了模型预测状态，而不能等同于真实液体不洒出的严格安全证明。硬防溢出约束和闭环液面反馈属于后续工作。

\subsection{预期结论}

本文预期证明的不是“\SPMPC{} 在所有导航指标上优于所有 planner”，而是以下更具体的结论：

\begin{quote}
  液体晃动不是普通轨迹平滑问题，而是具有动态记忆的状态预测问题。将低阶晃液状态嵌入滚动时域 \MPCC{} 后，移动底盘 local planner 能够在路径跟踪、路径进度、可执行控制、平滑性和预测液体响应之间进行权衡，从而在 smooth-only 之外进一步降低液体晃动。
\end{quote}

换言之，\SPMPC{} 不是“普通 local planner 加一个液体代价”，而是将 liquid-state prediction 放入在线局部规划层的晃液感知方法。

\subsection{后续工作}

后续工作可以围绕四个方向展开：

\begin{itemize}
  \item 引入真实液面反馈，实现 closed-loop liquid-state estimation 或参数自适应；
  \item 将 model-predicted slosh bound 扩展为更严格的防溢出约束和安全分析；
  \item 在 obstacle-aware、corridor/homotopy 或动态障碍物场景中扩展晃液感知局部规划；
  \item 评估不同液位、不同容器形状、多容器和真实实验室任务中的鲁棒性。
\end{itemize}
